\renewcommand\P{\mathbb{P}}

\begin{problem}[1]
  Suppose $X$ is a random variable with distribution function $\mu_X(t) = \frac{1}{\sqrt{2\pi}} e^{-t^2/2}$. Prove that $\E(X) = 0$ and $\E(X^2) = 1$.
\end{problem}

\begin{solution}
  Since
  \[%
    \mu_X(t) = \frac{1}{\sqrt{2\pi}} e^{-t^2/2},
  ,\]%
  this is the probability density function of the standard normal distribution. For a continuous random variable with density $f_X(t)$,
  \[%
    \E(X) = \int_{-\infty}^{\infty} t f_X(t) \dt \aand \E(X^2) = \int_{-\infty}^{\infty} t^2 f_X(t) \dt.
  .\]%
  For the first moment, we compute
  \[%
    \E(X) = \int_{-\infty}^{\infty} t \frac{1}{\sqrt{2\pi}} e^{-t^2/2} \dt = \frac{1}{\sqrt{2\pi}}\int_{-\infty}^{\infty} t e^{-t^2/2} \dt
  .\]%
  Observe that the function $te^{-t^2/2}$ is odd, since
  \[%
    (-t)e^{-(-t)^2/2} = -te^{-t^2/2}
  .\]%
  The integral of an odd function over a symmetric interval about $0$ is $0$, hence $\E(X) = 0$. For the second moment, we compute
  \[%
    \E(X^2) = \int_{-\infty}^{\infty} t^2 \frac{1}{\sqrt{2\pi}} e^{-t^2/2} \dt = \frac{1}{\sqrt{2\pi}}\int_{-\infty}^{\infty} t^2 e^{-t^2/2} \dt = \frac{\sqrt{2\pi}}{\sqrt{2\pi}} = 1
  .\]%
  Hence $\E(X) = 0$ and $\E(X^2) = 1$.
\end{solution}

\begin{problem}[2]
  Suppose $X$ is a random variable with law $\dd{\mu}_X = \sum_{k=0}^\infty \frac{\lambda^k}{k!} e^{-\lambda}\delta_k$, where $\delta_k$ is a Dirac mass at $k$ and $\lambda>0$.  Show that $\E(X) = \lambda$ and $\Var(X) = \lambda$. (\emph{Hint:} For the variance, first compute $\E\{X(X-1)\}$.) 
\end{problem}

\begin{solution}
  The law of $X$ is
  \[%
    \dd{\mu}_X = \sum_{k=0}^{\infty} \frac{\lambda^k}{k!}e^{-\lambda}\delta_k
  ,\]%
  so
  \[%
    \P(X = k) = \frac{\lambda^k}{k!}e^{-\lambda}
  ,\]%
  for $k = 0, 1, 2, \dots$. Thus $X$ has the Poisson distribution with parameter $\lambda$. For the expectation, by definition,
  \[%
    \E(X) = \sum_{k=0}^{\infty} k\,\P(X=k) = \sum_{k=0}^{\infty} k \frac{\lambda^k}{k!} e^{-\lambda}
  .\]%
  The $k = 0$ term vanishes. Observe that
  \[%
    k\frac{\lambda^k}{k!} = \lambda \frac{\lambda^{k-1}}{(k-1)!}
  .\]%
  Hence
  \[%
    \E(X) = \lambda e^{-\lambda}\sum_{k=1}^{\infty} \frac{\lambda^{k-1}}{(k-1)!}
  .\]%
  Let $j = k - 1$. Then
  \[%
    \sum_{k=1}^{\infty} \frac{\lambda^{k-1}}{(k-1)!} = \sum_{j=0}^{\infty} \frac{\lambda^j}{j!} = e^{\lambda}
  .\]%
  Therefore $\E(X) = \lambda$. For variance, we use the hint and compute $\E[X(X - 1)]$. By definition,
  \[%
    \E[X(X-1)] = \sum_{k=0}^{\infty} k(k-1)\frac{\lambda^k}{k!}e^{-\lambda}
  .\]%
  The terms for $k = 0, 1$ vanish, so
  \[%
    \E[X(X - 1)] = e^{-\lambda}\sum_{k=2}^{\infty} k(k - 1)\frac{\lambda^k}{k!}
  .\]%
  Observe that
  \[%
    k(k - 1)\frac{\lambda^k}{k!} = \lambda^2 \frac{\lambda^{k-2}}{(k-2)!}
  .\]%
  Thus
  \[%
    \E[X(X - 1)] = \lambda^2 e^{-\lambda}\sum_{k=2}^{\infty} \frac{\lambda^{k-2}}{(k-2)!}
  .\]%
  Let $j = k - 2$. Then
  \[%
    \sum_{k=2}^{\infty} \frac{\lambda^{k-2}}{(k - 2)!} = \sum_{j=0}^{\infty} \frac{\lambda^j}{j!} = e^{\lambda}
  .\]%
  Hence $\E[X(X - 1)] = \lambda^2$. Now observe that
  \[%
    X^2 = X(X - 1) + X
  .\]%
  Therefore
  \[%
    \E(X^2) = \E[X(X - 1)] + \E(X) = \lambda^2 + \lambda
  .\]%
  Finally,
  \[%
    \Var(X) = \E(X^2) - (\E(X))^2 = (\lambda^2 + \lambda) - \lambda^2 = \lambda
  .\]%
  Thus $\E(X) = \lambda$ and $\Var(X) = \lambda$.
\end{solution}

\begin{problem}[3]
  Suppose we are trying to determine the intensity of some source in the presence of background noise.  We therefore have two (independent) random variables $S$ and $N$ (the source and the noise), modeled as independent Poisson random variables with intensities $\lambda_s$ and $\lambda_n$, respectively.  By previous measurements, we have already determined $\lambda_n$.  Our goal is to determine $\lambda_s$.  Our sensor records only the total $X=S+N$. 
  \begin{enumerate}
    \item Show that $X$ is a Poisson random variable with intensity $\lambda_s+\lambda_n$. 
    \item Show that if the sensor records $N$ counts, then the maximum likelihood estimate for $\lambda_s$ is $N-\lambda_n$.  (Without loss of generality, you may assume $N>\lambda_n$.)  
  \end{enumerate}
  (To compute the maximum likelihood, you need to compute the probability that $X=N$ under the assumption that $\lambda_s=\lambda$, then optimize in the choice of $\lambda$.) 
\end{problem}

\begin{solution}[(i)]
  Let $S$ and $N$ be independent Poisson random variables with parameters $\lambda_s$ and $\lambda_n$, respectively. Thus
  \[%
    \P(S = k) = \frac{\lambda_s^k}{k!}e^{-\lambda_s} \aand \P(N = j) = \frac{\lambda_n^j}{j!}e^{-\lambda_n}
  .\]%
  Let $X = S + N$. For any integer $m \ge 0$,
  \[%
    \P(X = m) = \sum_{k=0}^{m} \P(S = k, N = m - k)
  .\]%
  Since $S$ and $N$ are independent,
  \[%
    \P(S = k, N = m - k) = \P(S = k)\P(N = m - k)
  .\]%
  Hence
  \[%
    \P(X = m) = \sum_{k=0}^{m} \frac{\lambda_s^k}{k!}e^{-\lambda_s} \frac{\lambda_n^{m-k}}{(m - k)!}e^{-\lambda_n}
  .\]%
  Factoring out the exponential term gives
  \[%
    \P(X = m) = e^{-(\lambda_s+\lambda_n)} \sum_{k=0}^{m} \frac{\lambda_s^k\lambda_n^{m-k}}{k!(m - k)!}
  .\]%
  Using
  \[%
    \binom{m}{k} = \frac{m!}{k!(m-k)!}
  ,\]%
  we rewrite the sum as
  \[%
    \sum_{k=0}^{m} \frac{\lambda_s^k\lambda_n^{m-k}}{k!(m - k)!} = \frac{1}{m!} \sum_{k=0}^{m} \binom{m}{k}\lambda_s^k\lambda_n^{m-k}
  .\]%
  Using the Binomial Theorem, we get
  \[%
    \P(X = m) = \frac{(\lambda_s + \lambda_n)^m}{m!}e^{-(\lambda_s + \lambda_n)}
  .\]%
  Thus $X$ is a Poisson random variable with intensity $\lambda_s + \lambda_n$.
\end{solution}

\begin{solution}[(ii)]
  Suppose the sensor records $N$ counts, so $X=N$. From part (i),
  \[%
    X \sim \text{Poisson}(\lambda_s + \lambda_n)
  .\]%
  To estimate $\lambda_s$, assume $\lambda_s=\lambda$ and compute the likelihood
  \[%
    L(\lambda) = \P(X = N \mid \lambda_s = \lambda)
  .\]%
  Using the Poisson distribution,
  \[%
    L(\lambda) = \frac{(\lambda + \lambda_n)^N}{N!}e^{-(\lambda + \lambda_n)}
  .\]%
  To maximize this likelihood, it is convenient to maximize the log-likelihood:
  \[%
    \ell(\lambda) = \log L(\lambda) = N\log(\lambda + \lambda_n) - (\lambda + \lambda_n) - \log(N!)
  .\]%
  Differentiating with respect to $\lambda$ gives
  \[%
    \ell'(\lambda) = \frac{N}{\lambda + \lambda_n}-1
  .\]%
  Setting this equal to zero yields
  \[%
    \frac{N}{\lambda + \lambda_n} = 1, \qquad \lambda + \lambda_n = N
  .\]%
  Thus $\lambda = N - \lambda_n$. Since we assume $N>\lambda_n$, this value is positive and therefore admissible. Hence the maximum likelihood estimate for $\lambda_s$ is $\hat{\lambda}_s = N - \lambda_n$.
\end{solution}

\begin{problem}[4]
  Let $f : \R^n \to \R$ be a continuously differentiable function. Suppose that $f$ has a unique global minimum at some point $x_* \in \R^n$. Let $x_0 \in \R^n$ be arbitrary and define $\{x_k\}$ recursively via
  \[%
    x_{k+1} = x_k - \lambda\nabla f(x_k)
  ,\]%
  where $\lambda > 0$.

  \begin{enumerate}
    \item Show that if the sequence $x_k$ converges, then it converges to $x_*$.
    \item Fix $n = 1$. Suppose in addition that $f$ is twice continuously differentiable and that $0 < m < f''(x) \leq M$ for some $m, M > 0$ and all $x \in \R$. Prove that the sequence $x_k$  defined above converges provided $\lambda$ is sufficiently small.
  \end{enumerate}

  (This is the gradient descent algorithm with learning rate $\lambda$.)
\end{problem}

\begin{solution}[(i)]
\end{solution}

\begin{solution}[(ii)]
\end{solution}

\begin{problem}[5]
  Let $f : \R^3 \to \R$ be infinitely differentiable and equal to zero outside of some ball. Show that
  \[%
    -\int_{\R^3} \frac{1}{|x|} \Delta f(x) \dx = 4\pi f(0)
  ,\]%
  where $\Delta$ is the Laplacian. This calculation shows that in three space dimensions, $-\Delta(1/|x|) = 4\pi\delta$ (where $\delta$ is the Dirac delta distribution).

  \noindent \emph{Hint.} Let $\epsilon > 0$ and split the integral into the set $|x| \leq \epsilon$ and the set $|x| > \epsilon$. On the set $|x| > \epsilon$, integrate by parts twice, using the identity
  \[%
    \int_S g \pd{j} f \dx = -\int_S f\pd{j} g \dx + \int_{\partial S} fg \nu_j \dA
  ,\]%
  where $\nu$ is the outward normal vector on $\partial S$ and $dA$ is surface measure.
\end{problem}

\begin{solution}
\end{solution}

\begin{problem}[6]
  Recall that the Green's function $G = G(t, x, y)$ solves
  \[%
    \frac{1}{c^2(x)}\pdv[2]{G}{t} - \Delta_x G = \delta(t)\delta(x - y)
  ,\]%
  with $G(t, x, y) = 0$ for $t < 0$. Show that 
  \[%
    u(t,x) = \iint G(t - s, x, y) n(s, y)\dy\ds
  ,\]%
  solves
  \[%
    \frac{1}{c^2(x)}\pdv[2]{u}{t} - \Delta_x u = n(t, x)
  .\]%
\end{problem}

\begin{solution}
  Let
  \[%
    u(t, x) = \iint G(t - s, x, y)n(s, y) \dd{y,s}
  \]%
  Define the operator
  \[%
    L = \frac{1}{c^2(x)}\pdv[2]{}{t} - \Delta_x
  .\]%
  Since differentiation with respect to $t$ and $x$ can be passed under the integral, we obtain
  \[%
    Lu(t, x) = \iint \left(\frac{1}{c^2(x)}\pdv[2]{}{t}G(t - s, x, y) - \Delta_x G(t - s, x, y)\right) n(s, y) \dd{y,s}
  .\]%
  Because $G$ satisfies
  \[%
    \frac{1}{c^2(x)}\pdv[2]{G}{t} - \Delta_x G = \delta(t)\delta(x - y)
  ,\]%
  we substitute $t - s$ for $t$ to obtain
  \[%
    \frac{1}{c^2(x)}\pdv[2]{}{t}G(t - s, x, y) - \Delta_x G(t - s, x, y) = \delta(t - s)\delta(x - y)
  .\]%
  Hence
  \[%
    Lu(t, x) = \iint \delta(t - s)\delta(x - y)n(s, y) \dd{y,s}
  \]%
  Using the defining property of the Dirac delta,
  \[%
    \int \delta(t - s)f(s) \ds = f(t) \aand \int \delta(x - y)g(y) \dy = g(x)
  ,\]%
  the integrals collapse, $Lu(t, x) = n(t, x)$ Thus
  \[%
    \frac{1}{c^2(x)}\pdv[2]{u}{t} - \Delta_x u = n(t, x)
  ,\]%
  so $u(t,x)$ solves the desired equation.
\end{solution}
